\begin{abstract}
The BlackSwanSuite benchmark evaluates vision-language models on abductive reasoning (inferring hidden causes) and defeasible reasoning (revising beliefs given new evidence) in videos of unexpected events.
While the benchmark is designed around \emph{video} understanding, we ask a more fundamental question: \textit{how much of BlackSwanSuite can be solved from text alone?}
We systematically evaluate five cognitively-motivated prompting strategies---Direct, Chain-of-Thought, Abductive Framing, Process-of-Elimination, and Counterfactual Contrastive---across three frontier language models (GPT-4o-mini, Claude~3.5~Haiku, Gemini~2.0~Flash) on the full 1{,}793-question MCQ validation set, \emph{without any visual input}.
Our text-only baselines achieve up to 55.5\% overall accuracy (vs.\ 33.3\% random chance), with Chain-of-Thought prompting yielding the strongest performance.
Strikingly, Reporter (defeasible) questions are substantially easier than Detective (abductive) questions in the text-only setting (69.0\% vs.\ 48.2\%), revealing that linguistic cues in answer options leak task-relevant information.
We further contribute a five-category error taxonomy---\emph{overweights prior}, \emph{hallucinated cause}, \emph{plausibility confusion}, \emph{fails to retract}, and \emph{overweights prior text}---and show that different prompting strategies induce qualitatively different error profiles.
Our results establish a rigorous text-only ceiling for BlackSwanSuite, quantify the ``visual gap'' that remains for multimodal models to exploit, and provide actionable guidance for future challenge participants.
All code and predictions are publicly released.
\end{abstract}
